\section{Automapper} 


% ===========================================
\subsection{checkPositionsAutomap}
% ===========================================

Validates the measurement plan generated by AutoMapper before moving the Suruga platform. It calculates and displays the absolute \texttt{x1} and \texttt{x2} positions for each device defined in the Excel file, but it does not connect to any instrument or move any axes.

\textbf{Source code:}
\href{https://github.com/niospinamev2/SurugaSeikiHex/blob/main/rutinas/checkPositionsAutomap.py}
{\texttt{checkPositionsAutomap.py}}.

\subsubsection{Workflow}

\begin{enumerate}
    \item Loads the chip layout and measurement plan from an Excel file using \texttt{automaper.load\_chip()}.

    \item Uses \texttt{origin\_x1} and \texttt{origin\_x2} as the reference origins.

    \item For each measurement, calculates the input and output offsets and the absolute positions of both arms.

    \item Displays the device, ports, and calculated positions, and waits for confirmation before displaying the next measurement.
\end{enumerate}

\subsubsection{Relevant Configuration}

%----------------------------------------------------------------
% Inicio tabla
%----------------------------------------------------------------

\footnotesize

\begin{longtblr}[
    caption = {Relevant configuration parameters for \texttt{checkPositionsAutomap}.},
    label = {tab:checkpositions-config}
]{
    colspec = {
        |Q[l, font=\ttfamily]|
        Q[l]|
        X[l]|
    },
    rowhead = 1,
    row{1} = {font=\bfseries},
    hlines,
}
%---------------------------
\SetRow{bg=gray9}
Parameter & Initial Value & Description \\
%---------------------------
\texttt{EXCEL\_FILE}
    & \texttt{chipMaps/SurugaTest.xlsx}
    & Chip layout and AutoMapper measurement plan. \\
%---------------------------
\texttt{origin\_x1} / \texttt{origin\_x2}
    & $0.0\,\mu\mathrm{m}$ / $0.0\,\mu\mathrm{m}$
    & Reference positions used to calculate the absolute coordinates.
\end{longtblr}

\normalsize

%----------------------------------------------------------------
% Fin tabla
%----------------------------------------------------------------

The absolute coordinates are calculated according to

\begin{equation}
    x_1 = \mathrm{origin}_{x1} - \mathrm{input}.x
\end{equation}

and

\begin{equation}
    x_2 = \mathrm{origin}_{x2} - \mathrm{output}.x.
\end{equation}

\subsubsection{Recommended Usage}

Run this routine whenever the Excel file, chip coupling, or origin convention changes. Compare the printed positions against the mechanical setup limits before running a routine that moves the machine.

\subsubsection{Results}

This routine does not create files or perform measurements. Its output consists of the position list displayed in the console.


% ===========================================
\subsection{moveAutomapOPM}
% ===========================================

Automatically iterates through the measurements defined by AutoMapper, positions the fibers at the input and output waveguides, performs alignment on both arms, and measures the average optical power using the optical power meter (OPM).

\textbf{Source code:}
\href{https://github.com/niospinamev2/SurugaSeikiHex/blob/main/rutinas/moveAutomapOPM.py}
{\texttt{moveAutomapOPM.py}}.

\subsection{Workflow}

\begin{enumerate}
    \item Loads the measurement plan from the AutoMapper Excel file.

    \item Connects to the Suruga platform and creates controllers for \texttt{x1}, \texttt{x2}, alignment, and the power meter.

    \item Loads \texttt{origin.txt} or uses the current position as the \texttt{WG0} origin.

    \item For each device, calculates the absolute positions of both arms and moves \texttt{x1} and \texttt{x2}.

    \item Performs a \textit{flat} alignment on the left arm followed by another alignment on the right arm.

    \item Reads ten samples from the OPM, calculates their average, and displays the result in the console.

    \item Requests confirmation before continuing and, once all measurements are completed, returns to \texttt{WG0}.
\end{enumerate}

\subsection{Relevant Configuration}

%----------------------------------------------------------------
% Inicio tabla
%----------------------------------------------------------------

\footnotesize

\begin{longtblr}[
    caption = {Relevant configuration parameters for \texttt{moveAutomapOPM}.},
    label = {tab:moveautomapopm-config}
]{
    colspec = {
        |Q[l, font=\ttfamily]|
        Q[l, font=\ttfamily]|
        X[l]|
    },
    rowhead = 1,
    row{1} = {font=\bfseries},
    hlines,
}

%---------------------------
\SetRow{bg=gray9}
Parameter & Initial Value & Description \\
%---------------------------
\texttt{EXCEL\_FILE}
    & \texttt{chipMaps/SurugaTest.xlsx}
    & Chip layout and measurement plan. \\
%---------------------------
\texttt{USE\_SAVED\_ORIGIN} / \texttt{SAVE\_ORIGIN}
    & \texttt{True} / \texttt{False}
    & Control whether \texttt{origin.txt} is loaded or saved. \\
%---------------------------
\texttt{ORIGIN\_FILE}
    & \texttt{origin.txt}
    & File containing the X1 and X2 coordinates of \texttt{WG0}. \\
%---------------------------
\texttt{flat\_left} / \texttt{flat\_right}
    & ---
    & Alignment parameters for each arm.
\end{longtblr}

\normalsize

%----------------------------------------------------------------
% Fin tabla
%----------------------------------------------------------------

\subsection{Results}

The average optical power for each device is printed to the console. The routine does not save the measurement values to a file.

\subsection{Precautions}

The routine moves both arms and requires manual confirmation between measurements. Verify the position plan using
\href{https://github.com/niospinamev2/SurugaSeikiHex/blob/main/rutinas/checkPositionsAutomap.py}{\texttt{checkPositionsAutomap}} before running it.


% ===========================================
\subsection{moveAutomapOPM\_OSA}
% ===========================================

Characterizes the devices defined by AutoMapper by combining Suruga platform
movement, automatic alignment, an average optical power measurement using the
OPM, and spectrum acquisition with the OSA.

\textbf{Source code:}
\href{https://github.com/niospinamev2/SurugaSeikiHex/blob/main/rutinas/moveAutomapOPM_OSA.py}
{\texttt{moveAutomapOPM\_OSA.py}}.

\subsection{Workflow}

\begin{enumerate}
    \item Loads the chip layout and measurement plan from Excel.

    \item Connects to the Suruga platform, OPM, and OSA, and configures the
    spectral sweep.

    \item Loads or saves the \texttt{WG0} origin in \texttt{origin2.txt}.

    \item For each device and repetition, positions \texttt{x1} and
    \texttt{x2} based on the corresponding input and output waveguides.

    \item Aligns the left and right arms.

    \item Runs an OSA sweep, retrieves the trace, and averages ten power
    readings.

    \item Saves the trace along with the device, ports, optical power, and
    repetition number as metadata.

    \item Returns to \texttt{WG0} and closes the OSA connection.
\end{enumerate}

\subsection{Relevant Configuration}

%----------------------------------------------------------------
% Inicio tabla
%----------------------------------------------------------------

\footnotesize

\begin{longtblr}[
    caption = {Relevant configuration parameters for
    \texttt{moveAutomapOPM\_OSA}.},
    label = {tab:moveautomapopm-osa-config}
]{
    colspec = {
        |Q[l, font=\ttfamily]|
        Q[l, font=\ttfamily]|
        X[l]|
    },
    rowhead = 1,
    row{1} = {font=\bfseries},
    hlines,
}

%---------------------------
\SetRow{bg=gray9}
Parameter & Initial Value & Description \\
%---------------------------
\texttt{DEBUG\_MODE}
    & \texttt{True}
    & Requests manual confirmations; when set to \texttt{False}, the routine
    runs automatically. \\
%---------------------------
\texttt{USE\_SAVED\_ORIGIN} / \texttt{SAVE\_ORIGIN}
    & \texttt{False} / \texttt{True}
    & Control whether the origin stored in \texttt{origin2.txt} is loaded or
    saved. \\
%---------------------------
\texttt{EXCEL\_FILE}
    & \texttt{chipMaps/SurugaTest01.xlsx}
    & Chip layout and AutoMapper measurement plan. \\
%---------------------------
\texttt{num\_repeticiones}
    & \texttt{1}
    & Number of complete passes through the measurement plan. \\
%---------------------------
\texttt{CENTER} / \texttt{SPAN}
    & 1550 nm / 90 nm
    & OSA spectral window. \\
%---------------------------
\texttt{RESOLUTION} / \texttt{SENSITIVITY}
    & 0.2 nm / -65 dBm
    & OSA resolution and sensitivity. \\
%---------------------------
\texttt{channel}
    & \texttt{1}
    & Power meter channel.
\end{longtblr}

\normalsize

%----------------------------------------------------------------
% Fin tabla
%----------------------------------------------------------------


\subsection{Results}

For each device, the routine saves a trace to:

\begin{center}
\texttt{data/<date>/medicion\_<repetition>\_<device>\_<input>\_<output>.txt}
\end{center}

The file contains the spectral data and metadata including the device, ports,
average optical power, and repetition number.

\subsection{Precautions}

With \texttt{DEBUG\_MODE=True}, the routine pauses before the OSA sweep,
between devices, and before returning to the origin.

Before using automatic mode, validate the movements with
\href{https://github.com/niospinamev2/SurugaSeikiHex/blob/main/rutinas/checkPositionsAutomap.py}
{\texttt{checkPositionsAutomap}}, verify the origin, and confirm that the OSA
VISA resource is available.